\section{Proofs}

\begin{proof}[Proof of Proposition~\ref{prop:gradual-automation}]
Equation~\eqref{eq:automation-odds} gives
\[
    \frac{s^M_t}{1-s^M_t}
    =
    \left(\frac{\nu}{1-\nu}\right)^{\sigma_{HM}}
    \left(\frac{w_t}{\xi}\right)^{\sigma_{HM}-1}.
\]
If $\sigma_{HM}>1$ and $w_t/\xi\to\infty$, the right-hand side diverges and
$s^M_t\to1$. Moreover, the unit cost in
\eqref{eq:research-unit-cost} converges to
\[
    P^E_t\longrightarrow
    \nu^{-\sigma_{HM}/(\sigma_{HM}-1)}\xi.
\]
The first expression in \eqref{eq:research-conditional-demands} then implies
$H_t/E_t\to0$. If $w_t/\xi$ remains bounded, the odds in
\eqref{eq:automation-odds} remain bounded, so $\sigma_{HM}>1$ does not by itself imply
$s^M_t\to1$.
\end{proof}

\begin{proof}[Proof of Proposition~\ref{prop:singularity-conditions}]
Let $z_t\equiv1/g_{A,t}$. Equation
\eqref{eq:ces-transition-growth-dynamics} implies the linear equation
\[
    \dot z_t=-\eta n z_t+D(s^M_t).
\]
Multiplying by $e^{\eta nt}$ and integrating gives
\eqref{eq:time-varying-singularity-condition}. If $D(s^M_t)\geq0$, the bracket
in that equation remains positive. If $D>0$ is constant, the logistic equation for
$g_A$ has the stable positive fixed point $\eta n/D$.

If $D=0$, direct integration gives
\[
    g_{A,t}=g_{A,0}e^{\eta nt},
    \qquad
    A_t
    =
    A_0\exp\left[
        \frac{g_{A,0}}{\eta n}
        \left(e^{\eta nt}-1\right)
    \right].
\]
Both expressions are finite for every finite $t$. If $D<0$ is constant, the
solution for the growth rate is
\[
    g_{A,t}
    =
    \frac{g_{A,0}e^{\eta nt}}
    {1-\frac{(-D)g_{A,0}}{\eta n}
        \left(e^{\eta nt}-1\right)}.
\]
The denominator reaches zero at the date in
\eqref{eq:constant-feedback-singularity-time}. Near that date, $g_A$ is
proportional to $(T^*-t)^{-1}$, so both $g_A$ and $A$ diverge. For a time-varying
$D(s^M_t)$, the same divergence occurs exactly when the bracket in
\eqref{eq:time-varying-singularity-condition} reaches zero.
\end{proof}

\begin{proof}[Proof of Proposition~\ref{prop:production-ces-regimes}]
Let $x_t\equiv X_t/L_{Y,t}$, $y_t\equiv Y_t/L_{Y,t}$, and write
$s_t\equiv s^X_t$. The monopoly condition can be written as
\[
    \frac{\xi x_t}{A_ty_t}
    =
    (1-\alpha)s_t[1-\varepsilon(s_t)],
    \qquad
    \varepsilon(s)
    =
    \frac{1-s}{\sigma_{XL}}+\alpha s.
\]

For $\sigma_{XL}<1$, revenue approaches zero both as $x_t\to0$ and as
$x_t\to\infty$. As marginal cost converges to zero, the monopolist approaches the
interior revenue-maximizing quantity, which satisfies $\varepsilon(\bar s^X)=1$.
Solving this condition gives
\[
    \bar s^X=\frac{1-\sigma_{XL}}{1-\alpha\sigma_{XL}}.
\]
Because the CES mapping between $s^X$ and $X/L_Y$ is one-to-one, the corresponding
service ratio is finite.

For $\sigma_{XL}=1$, the CES contribution $s^X_t$ equals the constant weight $\omega$.

For $\sigma_{XL}>1$, revenue is unbounded as $x_t\to\infty$ when marginal cost is zero,
so $x_t\to\infty$ and $s^X_t\to1$. Consequently,
$\varepsilon^X_t\to\alpha$, and the monopoly first-order condition implies
\[
    \frac{\xi U_t}{Y_t}\longrightarrow(1-\alpha)^2.
\]
Since
\[
    Z_t\sim
    \omega^{\sigma_{XL}/(\sigma_{XL}-1)}X_t,
\]
substituting $X_t=A_tU_t$ and the limiting compute share into production yields
\[
    \frac{Y_t}{K_t}
    \sim
    \left[
        \frac{\omega^{\sigma_{XL}/(\sigma_{XL}-1)}(1-\alpha)^2}{\xi}
    \right]^{(1-\alpha)/\alpha}
    A_t^{(1-\alpha)/\alpha}.
\]
This establishes \eqref{eq:high-sigma-ak-limit}.
\end{proof}

\begin{proof}[Derivation of Equation~\eqref{eq:bg-research-resource-share}]

In the Cobb--Douglas production benchmark, the envelope theorem and
\eqref{eq:cd-monopoly} imply
\[
    \pi^X_{A,t}=\frac{\beta^2Y_t}{A_t}.
\]
The research first-order condition in the AI-dominated limit gives
\[
    \xi=q_t\eta\frac{\dot A_t}{M_t},
    \qquad
    \frac{q_tA_t}{Y_t}
    =
    \frac{m_R^\infty}{\eta g_A^\infty}.
\]
Because the right-hand side is constant on a balanced-growth path,
$g_q=g_Y-g_A$. Dividing \eqref{eq:private-costate} by $q_t$, using
$r-\delta=\rho+\gamma^\infty$ and $g_Y=n+\gamma^\infty$, and substituting the
two expressions above gives
\[
    \rho-n
    =
    \frac{\beta^2\eta g_A^\infty}{m_R^\infty}
    -(1-\phi)g_A^\infty.
\]
Solving for $m_R^\infty$ yields
\eqref{eq:bg-research-resource-share}.
\end{proof}

\begin{proof}[Proof of Proposition~\ref{prop:knowledge-wedge}]
Evaluate the planner's marginal net benefit from research compute at the private
allocation. Using \eqref{eq:private-compute-foc},
\[
    Q_tF_{M,t}-\xi
    =
    (Q_t-q_t)F_{M,t}.
\]
If $Q_t>q_t$, the expression is positive and the planner increases research compute
locally relative to the private allocation.
\end{proof}
